\section*{Conclusions}

Media-only sampling of a minimal two-compartment liver-chip model identifies three combinations of four kinetic parameters, not the parameters themselves, and we give the explicit one-parameter family of equally consistent solutions. Independently measuring non-specific binding collapses this family to a point, making the remaining kinetic parameters structurally identifiable --- but, as we show directly by multi-start profile likelihood, this does not guarantee that they are practically estimable from realistic, finite, noisy data. Standard local diagnostics for practical identifiability (Fisher-information covariance approximations, single-start profile likelihood) proved unreliable near this structural degeneracy, and only a multi-start approach gave a trustworthy answer; we report this methodological finding as a caution for identifiability analyses of similar liver-chip and organ-on-a-chip models. We propose that liver-chip pharmacokinetic studies intended to support IVIVE report the raw data and loss-term characterization needed to assess structural identifiability, and perform a practical-identifiability analysis robust to the kind of degeneracy demonstrated here, rather than reporting point estimates alone. A preliminary finding --- that practical identifiability of $CL_{int}$ may depend strongly on the cell-compartment-to-media volume ratio, a fixed hardware property rather than a sampling choice --- is reported as an open question for future work.
